Une bobine circulaire de $N = \v{N}$ spires et de rayon \qty{\v{r}}{cm} a une résistance totale de \qty{\v{R}}{\ohm}. Ses extrémités sont reliées entre elles. Elle est placée dans un champ magnétique uniforme perpendiculaire à son plan, dont la grandeur passe de \qty{\v{B1}}{mT} à \qty{\v{B2}}{mT} en \qty{\v{dt}}{ms}, à un rythme constant.

\begin{enumerate}[label=\alph*)]
    \item Déterminez la variation du flux magnétique à travers une spire.
    \item Déterminez la force électromotrice induite dans la bobine.
    \item Déterminez le courant induit.
    \item Déterminez la charge totale qui a circulé dans la bobine pendant cette variation.
\end{enumerate}
